\documentclass[11pt]{article}
\newcommand{\DocShort}{Examen 1: conducción presurizada}
\usepackage{fontspec}
\setmainfont{Latin Modern Roman}
\setsansfont{Latin Modern Sans}
\usepackage[spanish,es-noshorthands]{babel}
\decimalpoint
\usepackage{geometry}
\geometry{letterpaper, top=2.3cm, bottom=2.3cm, left=2.5cm, right=2.5cm}
\usepackage[expansion=false]{microtype}
\usepackage{parskip}
\setlength{\parskip}{5pt}
\usepackage{amsmath,amssymb}
\usepackage{siunitx}
\usepackage{graphicx}
\usepackage{pdfpages}
\usepackage{pdflscape}
\usepackage{booktabs}
\usepackage{array}
\usepackage{longtable}
\usepackage{tabularx}
\usepackage{makecell}
\usepackage{caption}
\usepackage{enumitem}
\usepackage[dvipsnames,table]{xcolor}
\usepackage{fancyhdr}
\usepackage[hidelinks,pdfencoding=auto]{hyperref}
\usepackage{titlesec}
\usepackage{tcolorbox}
\tcbuselibrary{breakable}

\sisetup{
  group-separator = {,},
  group-minimum-digits = 4,
  output-decimal-marker = {.},
  per-mode = symbol
}

\definecolor{darkblue}{HTML}{16324F}
\definecolor{medblue}{HTML}{2C6C9B}
\definecolor{lightblue}{HTML}{EAF3FA}
\definecolor{warnbg}{HTML}{FFF3CD}
\definecolor{warnline}{HTML}{A25A00}
\definecolor{lightgrey}{HTML}{F7F7F7}
\definecolor{midgrey}{HTML}{666666}

\titleformat{\section}
  {\large\bfseries\color{darkblue}}
  {\thesection}{0.7em}{}[{\color{medblue}\titlerule[0.8pt]}]
\titleformat{\subsection}
  {\normalsize\bfseries\color{darkblue}}
  {\thesubsection}{0.7em}{}

\setlist[itemize]{topsep=2pt, itemsep=2pt, parsep=0pt}
\setlist[enumerate]{topsep=2pt, itemsep=2pt, parsep=0pt}

\pagestyle{fancy}
\fancyhf{}
\fancyhead[L]{\small\color{midgrey}ICV-441 Acueductos y Alcantarillados}
\fancyhead[R]{\small\color{midgrey}\DocShort}
\fancyfoot[C]{\small\color{midgrey}\thepage}
\setlength{\headheight}{14pt}
\renewcommand{\headrulewidth}{0.4pt}

\rowcolors{2}{lightgrey}{white}
\renewcommand{\arraystretch}{1.16}

\newtcolorbox{infobox}{
  breakable,
  colback=lightblue,
  colframe=medblue,
  boxrule=0.8pt,
  arc=3pt,
  left=8pt,
  right=8pt,
  top=7pt,
  bottom=7pt
}

\newtcolorbox{warnbox}{
  breakable,
  colback=warnbg,
  colframe=warnline,
  boxrule=0.8pt,
  arc=3pt,
  left=8pt,
  right=8pt,
  top=7pt,
  bottom=7pt
}

\newcommand{\doccover}[2]{
  \begin{center}
  {\small\color{midgrey}Pontificia Universidad Católica Madre y Maestra\\Escuela de Ingeniería Civil y Ambiental}

  \vspace{0.9cm}
  {\color{medblue}\rule{\linewidth}{2pt}}
  \vspace{0.25cm}

  {\Large\bfseries\color{darkblue}ICV-441 Acueductos y Alcantarillados}\\[0.35em]
  {\LARGE\bfseries\color{darkblue}#1}

  \vspace{0.25cm}
  {\color{medblue}\rule{\linewidth}{2pt}}
  \vspace{0.8cm}
  \end{center}

  \begin{center}
  \begin{tabular}{ll}
  \toprule
  \textbf{Asignatura} & ICV-441 Acueductos y Alcantarillados \\
  \textbf{Profesor} & Ricardo Rafael Hernández Moreira, PhD \\
  \textbf{Versión} & #2 \\
  \textbf{Entrega} & Un solo archivo PDF \\
  \bottomrule
  \end{tabular}
  \end{center}

  \vspace{0.35cm}
}

\newcommand{\introbox}[2]{
  \begin{infobox}
  \textbf{Propósito.} #1

  \vspace{0.5\baselineskip}

  \textbf{Qué se espera.} #2
  \end{infobox}
}

\AtBeginDocument{
  \hypersetup{
    pdftitle={\DocShort},
    pdfauthor={Ricardo Rafael Hernández Moreira},
    pdfsubject={ICV-441 Acueductos y Alcantarillados}
  }
}


\begin{document}

\doccover{Examen 1: diseño de conducción presurizada}{2026-07}
\introbox{Diseñar una conducción presurizada por gravedad, dada la geometría propuesta en planta y perfil, aplicando criterios normativos del INAPA, pérdidas mayores y menores, verificación de golpe de ariete y diseño selectivo de bloques de anclaje.}{Se espera una solución técnicamente consistente, económicamente razonable y completamente trazable, presentada en una memoria narrativa clara acompañada de una tabla principal de resultados y de un perfil longitudinal con línea piezométrica y línea de energía.}

\section{Objetivo}
Diseñar la tubería presurizada que conduce agua desde un reservorio con nivel de agua a cota \SI{193}{\meter} msnm hasta un tanque con nivel de agua a cota \SI{140}{\meter} msnm, usando el alineamiento y el perfil longitudinal suministrados como archivos adjuntos del examen.

\section{Normativa y jerarquía de referencias}
La referencia normativa principal y prevalente para este examen es el \emph{Reglamento Técnico para Diseño e Instalaciones Hidro-Sanitario del INAPA}, versión julio de 2018. En caso de ambigüedad, prevalecerá este reglamento sobre cualquier otra fuente.

Los estudiantes deberán citar en su memoria la sección o acápite y la página del reglamento o de la fuente técnica utilizada cuando un dato no aparezca explícitamente en esta consigna.

\begin{infobox}
\textbf{Referencias normativas mínimas a revisar en el reglamento INAPA.}
\begin{itemize}
\item Parámetros de diseño, velocidades permisibles, presiones mínimas y máximas: sección 6.3.2, pp. 78--80.
\item Procedimiento general para la determinación de consumos y gastos de cálculo: sección 6.3.6.1, p. 84.
\item Válvulas de corte y ventosas en líneas de conducción: pp. 70--72.
\end{itemize}
\end{infobox}

\section{Datos dados}
\subsection{Datos hidráulicos y físicos}
\begin{itemize}
\item Material de la tubería: acero.
\item Rugosidad absoluta del acero: $\varepsilon = \SI{0.15}{\milli\meter}$.
\item Temperatura del agua: \SI{20}{\celsius}.
\item Gravedad: $g = \SI{9.81}{\meter\per\second\squared}$.
\item Módulo volumétrico del agua: $E_{\text{agua}} = \SI{2.17e9}{\newton\per\meter\squared}$.
\item Módulo de elasticidad del acero: $E_{\text{acero}} = \SI{1.9e11}{\newton\per\meter\squared}$.
\item Coeficiente de Poisson del acero: $\mu_{\text{acero}} = 0.30$.
\item Modelo estructural para el análisis transitorio: tubería de pared delgada, anclada contra movimiento axial.
\end{itemize}

\subsection{Datos demográficos y de demanda}
\begin{itemize}
\item Población actual: \num{11800} habitantes.
\item Horizonte de diseño: \num{25} años.
\item Método de crecimiento poblacional: geométrico.
\item Tasa de crecimiento geométrico: \num{4}\,\% anual.
\item Tipo de población: urbana.
\item Se aplicarán los factores de día máximo y hora máxima del Reglamento del INAPA.
\item Se supondrá un \SI{20}{\percent} de fugas.
\item Solo se considerará uso doméstico.
\item La conducción se diseñará hidráulicamente para el \textbf{caudal máximo diario}.
\end{itemize}

\subsection{Datos geométricos}
\begin{itemize}
\item Nivel de agua en el reservorio: \SI{193}{\meter} msnm.
\item Nivel de agua en el tanque final: \SI{140}{\meter} msnm.
\item La progresiva \texttt{0+000} corresponde al reservorio.
\item Las progresivas avanzan en sentido aguas abajo.
\item El alineamiento y el perfil longitudinal forman parte de los archivos adjuntos del examen y deben citarse, pero no reproducirse dentro de la memoria como apéndice de datos brutos.
\end{itemize}

\section{Método hidráulico obligatorio}
\begin{enumerate}
\item La pérdida por fricción se calculará con la ecuación de Darcy-Weisbach.
\item El factor de fricción de Darcy se determinará con la ecuación de Colebrook-White.
\item Alternativamente, podrá usarse la ecuación explícita de Swamee-Jain o el diagrama de Moody, siempre que el procedimiento quede claramente explicado.
\item El uso de la ecuación de Hazen-Williams será penalizado con \textbf{el 100\% de la calificación total del examen}.
\end{enumerate}

\section{Criterios de diseño y aceptación}
\subsection{Velocidades, presiones y balance de energía}
\begin{itemize}
\item Las velocidades permisibles y las presiones mínimas y máximas deberán cumplir con el Reglamento del INAPA.
\item Para zona urbana, se tomará como referencia la recomendación del reglamento de diseñar con velocidad máxima de \SI{1.60}{\meter\per\second}, salvo justificación normativa mejor sustentada.
\item Debe verificarse que la línea piezométrica permanezca por encima de la tubería en todo el trazado.
\item Debe presentarse un perfil longitudinal con la tubería, la línea piezométrica y la línea de energía.
\item La diferencia entre la carga disponible y la suma de pérdidas totales del sistema no podrá exceder \SI{0.10}{\meter}.
\item Las presiones calculadas deberán ser menores que la presión admisible de la clase de tubería seleccionada. Para el análisis transitorio y el cálculo del coeficiente $k$ de interacción fluido-estructura, véase el Cuadro~\ref{tab:k-golpe-ariete}, \emph{Valores de $k$ tomados de mecánica de sólidos para el análisis de golpe de ariete}; esta referencia se cita aquí y también más adelante en la sección de golpe de ariete para evitar ambigüedad de numeración.
\end{itemize}

\subsection{Diámetros y división de tramos}
\begin{itemize}
\item Se utilizarán diámetros comerciales disponibles para acero, según el Cuadro~\ref{tab:nps-dn}, \emph{Diámetros comerciales de acero para referencia del examen}.
\item El diseño deberá hacerse con diámetro interior.
\item Se permite más de un cambio de diámetro, pero se recomienda procurar una solución con un solo cambio cuando sea técnica y económicamente razonable.
\item Para la tabla principal, un nuevo tramo deberá iniciarse cada vez que ocurra cualquiera de los siguientes eventos: cambio de diámetro, accesorio con pérdida menor, quiebre del perfil, cambio de dirección en planta, ubicación de válvula, ventosa, purga o cualquier otro punto singular relevante.
\item Las piezas especiales entre estaciones deberán reportarse con su progresiva.
\end{itemize}

\subsection{Pérdidas menores y accesorios}
\begin{itemize}
\item Deben considerarse pérdidas menores.
\item No se permite el uso de longitudes equivalentes.
\item Las pérdidas menores deberán calcularse mediante coeficientes $K$.
\item Deben considerarse todos los codos sugeridos por el archivo de alineamiento.
\item Para los codos, utilizar $K=\zeta$ de Hinds según la Figura~\ref{fig:hinds-usbr}, \emph{Anexo de referencia para coeficientes de pérdida por codos según Hinds/USBR}.
\item La entrada a la tubería desde el reservorio se asumirá recta.
\item La entrada al tanque se considerará de borde recto.
\item El tipo de válvula de corte será de compuerta.
\item La disposición de válvulas, ventosas y otros accesorios deberá justificarse con base en el Reglamento del INAPA.
\end{itemize}

\section{Análisis de golpe de ariete}
\begin{itemize}
\item Debe realizarse un análisis de golpe de ariete considerando solo la tubería de mayor diámetro y la longitud correspondiente a esa tubería.
\item Se exige evaluar un cierre instantáneo usando únicamente la ecuación de Joukowsky:
\[
\Delta p = \rho V a
\]
donde $\rho$ es la densidad del agua, $V$ es la velocidad media en la tubería y $a$ es la celeridad de la onda.
\item Debe proponerse además una recomendación de cierre gradual tal que reduzca la sobrepresión y cumpla simultáneamente dos condiciones: que el esfuerzo en la tubería no exceda el esfuerzo admisible y que el tiempo de cierre sea mayor que el tiempo crítico.
\item Para esa recomendación, deberá compararse el tiempo de cierre propuesto con el tiempo crítico $2L/a$.
\item La selección de la clase de tubería se hará con el criterio:
\[
P_{\text{estática}} + \Delta P_{\text{transitoria}} + \text{margen}
\leq P_{\text{admisible de la clase}}
\]
\item Para el análisis de golpe de ariete deberá usarse el espesor correspondiente a la clase seleccionada del Cuadro~\ref{tab:clases-acero}, \emph{Cuadro simplificado de clase de tubería de acero para el examen}. Para el coeficiente $k$ aplicable a la condición de soporte adoptada, use el Cuadro~\ref{tab:k-golpe-ariete}.
\end{itemize}

\section{Bloques de anclaje}
\begin{itemize}
\item Deben diseñarse bloques de anclaje para los tres cambios de dirección de mayor ángulo.
\item El objetivo es determinar el área del bloque contra la zanja.
\item Se adoptará un factor de seguridad de \num{1.2}.
\item Use como suelo de diseño una mezcla densa de arena y grava con $\phi = 40^\circ$.
\item A \SI{1}{\meter} de profundidad, use $\sigma_a = \num{1800}\ \text{lb/ft}^2 = \SI{0.0862}{\mega\pascal}$.
\item Para este examen se pide el modelo simplificado de empuje por presión interna equilibrado por la presión admisible del terreno sobre el área normal a la zanja, usando como referencia el Cuadro~\ref{tab:presion-suelo}, \emph{Presión admisible del suelo traducida a unidades SI}.
\end{itemize}

\section{Memoria de cálculo y formato de entrega}
\begin{itemize}
\item El entregable será un único archivo PDF.
\item El nombre del archivo deberá ser exactamente:
\[
\texttt{1920\_ICV\_441\_\_Examen\_1\_[Nombre\_Apellido].pdf}
\]
\item Solo se aceptarán archivos PDF.
\item El formato de papel será carta.
\item El tamaño de letra deberá estar entre \num{10} y \num{12} puntos.
\item La memoria deberá ser económica, pero clara.
\item Todas las variables deberán definirse.
\item Todas las unidades deberán reportarse.
\item Se espera un máximo de tres decimales en los resultados.
\item Para valores menores que uno, cuando resulte conveniente, puede usarse notación científica.
\item Se alienta el uso correcto de prefijos SI.
\item Deben mostrarse las fórmulas utilizadas.
\item La tabla de resultados deberá ocupar una posición predominante dentro de la estructura de la memoria.
\item Se permite usar Excel, pero sus resultados deberán contextualizarse dentro de una memoria narrativa.
\item La entrega de una impresión de Excel como sustituto de la memoria narrativa será penalizada con \textbf{el 100\% de la calificación total del examen}.
\item Se permite usar EPANET para verificar cómputos.
\item Se permite usar Civil 3D o AutoCAD para la elaboración de dibujos.
\end{itemize}

\section{Actividades requeridas}
\begin{enumerate}
\item Determinar la población futura con el método geométrico y el horizonte de diseño dado.
\item Determinar el caudal medio diario, el caudal máximo diario y el caudal máximo horario con base en el Reglamento del INAPA.
\item Seleccionar una combinación adecuada de diámetros y longitudes de tubería tal que la carga disponible sea balanceada por las pérdidas mayores y menores totales del sistema.
\item Verificar velocidades y presiones, incluyendo ausencia de subpresiones.
\item Construir el perfil longitudinal con la tubería, la línea piezométrica y la línea de energía.
\item Seleccionar la clase de tubería considerando presión estática, golpe de ariete y margen.
\item Diseñar los tres bloques de anclaje requeridos.
\end{enumerate}

\section{Rúbrica de evaluación}
\begin{center}
\begin{tabularx}{\textwidth}{>{\raggedright\arraybackslash}p{0.49\textwidth} >{\raggedleft\arraybackslash}p{0.18\textwidth} >{\raggedright\arraybackslash}X}
\toprule
\textbf{Componente} & \textbf{Peso} & \textbf{Alcance} \\
\midrule
Análisis poblacional y determinación de caudales & 20\% & Incluye población futura, caudal medio, máximo diario y máximo horario. \\
Determinación de diámetros & 30\% & Selección de diámetros con base en el balance hidráulico y en la normativa aplicable. \\
Determinación de longitudes por diámetro & 20\% & Ajuste de longitudes para balance energético y trazabilidad por tramos. \\
Determinación de clase de tubería & 15\% & Incluye presión estática, transitoria y selección final de clase. \\
Diseño de bloques de anclaje & 15\% & Incluye empuje, presión admisible del suelo y área requerida. \\
\bottomrule
\end{tabularx}
\end{center}

\section{Errores críticos}
\begin{warnbox}
Los siguientes errores se consideran críticos:
\begin{itemize}
\item omitir pérdidas menores;
\item no revisar golpe de ariete;
\item no justificar cambios de diámetro;
\item no especificar clase de tubería;
\item no presentar el gráfico con línea piezométrica y línea de energía;
\item no diseñar los anclajes requeridos.
\end{itemize}

Además, las siguientes acciones penalizan con \textbf{el 100\% de la calificación total del examen}:
\begin{itemize}
\item usar la ecuación de Hazen-Williams en lugar del método exigido;
\item entregar una impresión de Excel como sustituto de la memoria narrativa.
\end{itemize}
\end{warnbox}

\section{Referencias externas permitidas}
Todo dato no dado explícitamente en esta consigna podrá obtenerse de tablas o libros técnicos, siempre que la fuente se cite correctamente. Sin embargo, para los datos incluidos en los anexos de este examen deberán usarse exclusivamente los valores aquí suministrados.

\clearpage
\tableofcontents
\clearpage
\clearpage
\section*{Anexo A. Diámetros comerciales de acero para el examen}
\addcontentsline{toc}{section}{Anexo A. Diámetros comerciales de acero para el examen}

\begin{infobox}
Para fines de este examen, el diámetro comercial se identificará con el \emph{nominal pipe size} (NPS) y su equivalente aproximado en diámetro nominal (DN). El diseño hidráulico debe realizarse siempre con \textbf{diámetro interior}, no con diámetro nominal ni con diámetro exterior.
\end{infobox}

\begin{table}[htbp]
\centering
\small
\caption{Diámetros comerciales de acero para referencia del examen}
\label{tab:nps-dn}
\begin{tabular}{ccc}
\toprule
\textbf{NPS (in)} & \textbf{DN aproximado} & \textbf{Diámetro exterior (mm)} \\
\midrule
3 & 80 & 88.9 \\
4 & 100 & 114.3 \\
5 & 125 & 141.3 \\
6 & 150 & 168.3 \\
8 & 200 & 219.1 \\
10 & 250 & 273.0 \\
12 & 300 & 323.9 \\
14 & 350 & 355.6 \\
16 & 400 & 406.4 \\
18 & 450 & 457.2 \\
20 & 500 & 508.0 \\
22 & 550 & 558.8 \\
24 & 600 & 609.6 \\
26 & 650 & 660.4 \\
28 & 700 & 711.2 \\
30 & 750 & 762.0 \\
32 & 800 & 812.8 \\
34 & 850 & 863.6 \\
36 & 900 & 914.4 \\
\bottomrule
\end{tabular}
\end{table}

\clearpage
\section*{Anexo B. Cuadro simplificado de clase de tubería de acero para el examen}
\addcontentsline{toc}{section}{Anexo B. Cuadro simplificado de clase de tubería de acero para el examen}

\begin{warnbox}
Este cuadro corresponde al \textbf{Cuadro~2, Cuadro simplificado de clase de tubería de acero para el examen}, y se provee exclusivamente para este examen. Para seleccionar la tubería, use solo los valores de esta tabla. No sustituya estos datos por catálogos externos. La presión admisible de la clase seleccionada debe ser mayor o igual que la suma de la presión estática, la sobrepresión transitoria y el margen adoptado por el diseñador.
\end{warnbox}

\begin{table}[htbp]
\centering
\footnotesize
\caption{Cuadro simplificado de clase de tubería de acero para el examen}
\label{tab:clases-acero}
\begin{tabular}{cccccc|cccccc}
\toprule
\makecell{\textbf{NPS}\\\textbf{(in)}} &
\textbf{DN} &
\textbf{Clase} &
\makecell{\textbf{$t$}\\\textbf{(mm)}} &
\makecell{\textbf{$D_i$}\\\textbf{(mm)}} &
\makecell{\textbf{$P_{\mathrm{adm}}$}\\\textbf{(MPa)}} &
\makecell{\textbf{NPS}\\\textbf{(in)}} &
\textbf{DN} &
\textbf{Clase} &
\makecell{\textbf{$t$}\\\textbf{(mm)}} &
\makecell{\textbf{$D_i$}\\\textbf{(mm)}} &
\makecell{\textbf{$P_{\mathrm{adm}}$}\\\textbf{(MPa)}} \\
\midrule
3  & 80  & 150 & 4.78  & 79.34  & 1.03 & 16 & 400 & 150 & 7.92  & 390.56 & 1.03 \\
3  & 80  & 200 & 5.56  & 77.78  & 1.38 & 16 & 400 & 200 & 9.53  & 387.34 & 1.38 \\
3  & 80  & 250 & 6.35  & 76.20  & 1.72 & 16 & 400 & 250 & 11.11 & 384.18 & 1.72 \\
4  & 100 & 150 & 4.78  & 104.74 & 1.03 & 18 & 450 & 150 & 7.92  & 441.36 & 1.03 \\
4  & 100 & 200 & 5.56  & 103.18 & 1.38 & 18 & 450 & 200 & 9.53  & 438.14 & 1.38 \\
4  & 100 & 250 & 6.35  & 101.60 & 1.72 & 18 & 450 & 250 & 11.11 & 434.98 & 1.72 \\
6  & 150 & 150 & 4.78  & 158.74 & 1.03 & 20 & 500 & 150 & 9.53  & 488.94 & 1.03 \\
6  & 150 & 200 & 5.56  & 157.18 & 1.38 & 20 & 500 & 200 & 11.11 & 485.78 & 1.38 \\
6  & 150 & 250 & 6.35  & 155.60 & 1.72 & 20 & 500 & 250 & 12.70 & 482.60 & 1.72 \\
8  & 200 & 150 & 6.35  & 206.40 & 1.03 & 24 & 600 & 150 & 9.53  & 590.54 & 1.03 \\
8  & 200 & 200 & 7.92  & 203.26 & 1.38 & 24 & 600 & 200 & 11.11 & 587.38 & 1.38 \\
8  & 200 & 250 & 9.53  & 200.04 & 1.72 & 24 & 600 & 250 & 12.70 & 584.20 & 1.72 \\
10 & 250 & 150 & 6.35  & 260.30 & 1.03 & 30 & 750 & 150 & 11.11 & 739.78 & 1.03 \\
10 & 250 & 200 & 7.92  & 257.16 & 1.38 & 30 & 750 & 200 & 12.70 & 736.60 & 1.38 \\
10 & 250 & 250 & 9.53  & 253.94 & 1.72 & 30 & 750 & 250 & 14.27 & 733.46 & 1.72 \\
12 & 300 & 150 & 6.35  & 311.20 & 1.03 & 36 & 900 & 150 & 12.70 & 889.00 & 1.03 \\
12 & 300 & 200 & 7.92  & 308.06 & 1.38 & 36 & 900 & 200 & 14.27 & 885.86 & 1.38 \\
12 & 300 & 250 & 9.53  & 304.84 & 1.72 & 36 & 900 & 250 & 15.88 & 882.64 & 1.72 \\
14 & 350 & 150 & 7.92  & 339.76 & 1.03 &    &     &     &       &        &      \\
14 & 350 & 200 & 9.53  & 336.54 & 1.38 &    &     &     &       &        &      \\
14 & 350 & 250 & 11.11 & 333.38 & 1.72 &    &     &     &       &        &      \\
\bottomrule
\end{tabular}
\caption*{\footnotesize Donde $t$ es el espesor mínimo, $D_i$ es el diámetro interior de diseño y $P_{\mathrm{adm}}$ es la presión admisible de la clase.}
\end{table}

\clearpage
\section*{Anexo C. Presión admisible del suelo traducida a SI}
\addcontentsline{toc}{section}{Anexo C. Presión admisible del suelo traducida a SI}

\begin{infobox}
Las presiones admisibles de suelo de referencia para el diseño de bloques de anclaje se presentan en el \textbf{Cuadro~3}. Para este examen, use como caso base la fila de \textbf{mezcla densa de arena y grava} a \SI{1}{\meter} de profundidad, con $\phi = 40^\circ$ y $\sigma_a = \SI{0.0862}{\mega\pascal}$.
\end{infobox}

\begin{table}[htbp]
\centering
\footnotesize
\caption{Presión admisible del suelo traducida a unidades SI}
\label{tab:presion-suelo}
\begin{tabular}{lcccc}
\toprule
\textbf{Material natural del suelo} & \textbf{0.61 m} & \textbf{0.91 m} & \textbf{1.22 m} & \textbf{1.52 m} \\
 & \multicolumn{4}{c}{\textbf{Presión admisible (kPa)}} \\
\midrule
Roca sana & 383.1 & 478.8 & 478.8 & 478.8 \\
Mezcla densa de arena y grava ($\phi = 40^\circ$) & 57.4 & 86.2 & 114.9 & 143.6 \\
Arena densa fina a gruesa ($\phi = 35^\circ$) & 38.3 & 57.4 & 79.0 & 100.6 \\
Mezcla de limo y arcilla ($\phi = 25^\circ$) & 23.9 & 33.5 & 45.5 & 57.4 \\
Arcilla blanda y suelos orgánicos ($\phi = 10^\circ$) & 9.6 & 14.4 & 19.2 & 23.9 \\
\bottomrule
\end{tabular}
\end{table}

\clearpage
\section*{Anexo D. Valores de k para análisis de golpe de ariete}
\addcontentsline{toc}{section}{Anexo D. Valores de k para análisis de golpe de ariete}

\begin{table}[htbp]
\centering
\footnotesize
\caption{Valores de $k$ tomados de mecánica de sólidos para el análisis de golpe de ariete}
\label{tab:k-golpe-ariete}
\begin{tabularx}{\textwidth}{>{\raggedright\arraybackslash}p{0.38\textwidth} >{\centering\arraybackslash}X >{\centering\arraybackslash}p{0.18\textwidth}}
\toprule
\textbf{Condiciones de apoyo} & \textbf{$k$ (paredes gruesas)} & \textbf{$k$ (paredes delgadas)} \\
\midrule
Anclada solo en el extremo aguas arriba. Libre para moverse longitudinalmente &
$\dfrac{2t}{D}(1+\mu)+\left(\dfrac{D}{D+t}\right)-\left(1-\dfrac{\mu}{2}\right)$ &
$1-\dfrac{\mu}{2}$ \\
Anclada contra movimiento axial (longitudinal), de modo que el esfuerzo axial sea cero &
$\dfrac{2t}{D}(1+\mu)+\left(\dfrac{D}{D+t}\right)(1-\mu^2)$ &
$1-\mu^2$ \\
Con juntas de expansión en toda la longitud (esfuerzo axial = 0) &
$\dfrac{2t}{D}(1+\mu)+\left(\dfrac{D}{D+t}\right)(1)$ &
$1$ \\
En túneles ($t=\infty$) &
$\dfrac{2t}{D}(1+\mu)$ &
No aplica \\
\bottomrule
\end{tabularx}
\caption*{\footnotesize Nota: $\dfrac{D}{D+t}= \dfrac{1}{1+t/D}$; $\mu$ es la razón de Poisson y típicamente toma valores entre 0.25 y 0.30 para muchos materiales de tubería.}
\end{table}

\clearpage
\section*{Anexo E. Coeficientes de pérdida por codos según Hinds/USBR}
\addcontentsline{toc}{section}{Anexo E. Coeficientes de pérdida por codos según Hinds/USBR}

\begin{figure}[h]
\centering
\rotatebox{90}{\includegraphics[page=1,scale=0.73, trim={2.5cm 1.8cm 2.5cm 2cm},clip]{../insumos/Head loss coefficients for pipe bends.pdf}}
\caption{Anexo de referencia para coeficientes de pérdida por codos según Hinds/USBR, citado en la sección de pérdidas menores y accesorios.}
\label{fig:hinds-usbr}
\end{figure}

\clearpage
\section*{Anexo F. Encabezados mínimos de la tabla entregable}
\addcontentsline{toc}{section}{Anexo F. Encabezados mínimos de la tabla entregable}

\begin{longtable}{p{0.95\textwidth}}
\toprule
\textbf{Encabezados sugeridos para la tabla principal del entregable} \\
\midrule
Ítem \\
Progresiva inicial \\
Progresiva final \\
Longitud del tramo (\si{\meter}) \\
Tipo de elemento \\
Diámetro nominal \\
Diámetro interior de diseño (\si{\milli\meter}) \\
Material \\
Clase de tubería \\
Caudal (\si{\liter\per\second} o \si{\meter\cubed\per\second}) \\
Velocidad (\si{\meter\per\second}) \\
Pérdida por fricción (\si{\meter}) \\
Pérdida menor (\si{\meter}) \\
Pérdida total acumulada (\si{\meter}) \\
Cota del terreno (\si{\meter}) \\
Cota de la tubería (\si{\meter}) \\
Carga piezométrica (\si{\meter}) \\
Carga total o línea de energía (\si{\meter}) \\
Presión (\si{\meter} de columna de agua o \si{\kilo\pascal}) \\
Pieza especial / válvula / ventosa / purga \\
Observaciones \\
\bottomrule
\end{longtable}


\end{document}
